\subsection{谱投影}

设 \(u\) 是复希尔伯特空间 E 上的正规部分算子。设 A 是 \(\mathrm{Sp}(u)\) 的普遍可测子集，\(\varphi_{\mathrm{A}}\) 是它的特征函数。由于 \(\varphi_{\mathrm{A}}\) 有界且满足 \(\varphi_{\mathrm{A}}^{2} = \varphi_{\mathrm{A}}\) ，E 的自同态 \(\varphi_{\mathrm{A}}(u)\) 是 E 的正交投影。它称为由 A 定义的 \(u\) 的谱投影。我们记 \(p_{\mathrm{A}} = \varphi_{\mathrm{A}}(u)\) 。我们有 \(p_{\emptyset} = 0\) 及 \(p_{\mathrm{Sp}(u)} = 1_{\mathrm{E}}\) 。

设 A 是 C 的普遍可测子集。由 A 定义的 u 的谱投影是谱投影 \(p_{\mathrm{Sp}(u)\cap\mathrm{A}}\) 。它也简单地记作 \(p_{A}\) 。对每个函数 \(f \in \mathcal{L}_{\mathrm{u}}(\mathrm{Sp}(u))\) ，对所有 \(x \in E\) 及每个 \(y \in \mathrm{dom}(f(u))\) ，我们有公式

\[
\langle p _ {\mathrm{A}} (x) \mid f (u) y \rangle = \int_ {\operatorname{Sp} (u) \cap \mathrm{A}} f d \mu ,\tag{10}
\]

其中 \(\mu\) 是 \((x,y)\) 关于 \(u\) 的谱测度（p. 271 的公式 (7)）。设 A 与 B 是 \(\mathrm{Sp}(u)\) 的普遍可测子集。由于 \(\varphi_{\mathrm{A}}\varphi_{\mathrm{B}} = \varphi_{\mathrm{A}\cap \mathrm{B}}\) ，我们有 \(p_{\mathrm{A}} \circ p_{\mathrm{B}} = p_{\mathrm{A}\cap \mathrm{B}}\) （IV, p. 275 的命题 5, \(c\) )）。特别地，若 A 与 B 不相交，则 \(p_{\mathrm{A}}\) 的像与 \(p_{\mathrm{B}}\) 的像正交。

命题 8. --- 设 \((\mathrm{A}_{i})_{i\in\mathrm{I}}\) 是 \(\mathrm{Sp}(u)\) 的两两不相交的普遍可测子集组成的可数族，并设 \(p_{i}\) 是由 \(A_{i}\) 定义的 u 的谱投影。诸集合 \(A_{i}\) 的并 A 是 \(\mathrm{Sp}(u)\) 的普遍可测子集，且级数 \(\sum_{i}p_{i}\) 收敛于 \(p_{A}\) ，这里的 \(\mathcal{L}(\mathrm{E})\) 赋以逐点收敛拓扑。

由 INT, IV, p. 177, § 5, 第 4 目, 推论 2，集合 A 是普遍可测的。级数 \(\sum_{i} \varphi_{\mathrm{A}_{i}}\) 逐点收敛于 \(\varphi_{\mathrm{A}}\) ，且其部分和以 1 为界。因此断言由 IV, p. 276 的命题 6 得出。

命题 9. --- 设 A 是 Sp(u) 的闭子集。设 \(\varphi_{\mathrm{A}}\) 是 A 的特征函数，\(p_{\mathrm{A}} = \varphi_{\mathrm{A}}(u)\) 是由 A 定义的 u 的谱投影。用 \(\mathrm{E}_{\mathrm{A}}\) 表示 \(p_{\mathrm{A}}\) 的像。它是 E 的闭子空间。

\begin{enumerate}
\def\labelenumi{\alph{enumi})}
\item
  子空间 \(E_{A}\) 是使 x 关于 u 的谱测度的支撑包含在 A 中的那些 \(x \in E\) 所成的集；
\item
  若 A 在 \(\mathbf C\) 中有界，则 \(E_{A}\) 包含在 u 的定义域中；
\item
  对属于 u 的定义域的所有 x，我们有 \(p_{\mathrm{A}}(x) \in \operatorname{dom}(u)\) 且 \(u(p_{\mathrm{A}}(x)) \in \mathrm{E}_{\mathrm{A}}\) ；特别地，\(u(x) \in \mathrm{E}_{\mathrm{A}}\) ，只要 \(x \in \operatorname{dom}(u) \cap \mathrm{E}_{\mathrm{A}}\) 。
\end{enumerate}

对 E 中的 \(x\) ，我们用 \(\mu_x\) 表示 \(x\) 关于 \(u\) 的谱测度。

我们来证明 \(a\) )。设 \(x\) 是 \(\mathrm{E}_{\mathrm{A}}\) 的元素。设 \(z \in \mathbf{C} - \mathbf{A}\) ，并设 U 是 \(z\) 的与 A 不相交的相对紧开邻域。对每个支撑包含在 U 中的函数 \(f \in \mathcal{K}(\mathbf{C})\) ，我们有

\[
\int_ {\mathrm{Sp} (u)} f \mu_ {x} = \langle x | f (u) x \rangle = \langle p _ {\mathrm{A}} (x) | f (u) x \rangle = \int_ {\mathrm{Sp} (u) \cap \mathrm{A}} f \mu_ {x} = 0
\]

（依据公式 (10)）。这表明 \(z\) 不属于 \(\mu_x\) 的支撑。因此，\(\mu_x\) 的支撑包含在 A 中。

反之，设 \(x \in E\) 使 \(\mu_{x}\) 的支撑包含在 A 中。由此可得

\[
\langle x \mid p _ {\mathrm{A}} (x) \rangle = \int_ {\operatorname{Sp} (u) \cap \mathrm{A}} \mu_ {x} = \int_ {\operatorname{Sp} (u)} \mu_ {x} = \langle x \mid x \rangle ,
\]

\((loc. \ cit.) \ \text{因此} \ \|p_{\mathrm{A}}(x)-x\|^{2} = \|p_{\mathrm{A}}(x)\|^{2} - \|x\|^{2} \leqslant 0 \ \text{从而 } p_{\mathrm{A}}(x) = x, \ \text{即 } x \in \mathrm{E}_{\mathrm{A}}.\)

断言 b) 由 a) 及 IV, p. 273 的注得出。

我们来证明 \(c\) )。\(u\) 的定义域是这样的 \(x \in \mathrm{E}\) 所成的集：\(\operatorname{Sp}(u)\) 的恒等函数属于 \(\mathcal{L}^2(\operatorname{Sp}(u), \mu_x)\) （loc. cit.）。由于 \(\mu_{p_{\mathrm{A}}(x)} = \varphi_{\mathrm{A}} \cdot \mu_x\) 对 \(x \in \mathrm{E}\) 成立（IV, p. 275 的命题 5 的推论），我们有 \(p_{\mathrm{A}}(x) \in \operatorname{dom}(u)\) ，只要 \(x \in \operatorname{dom}(u)\) 。

设 \(x \in \operatorname{dom}(u)\) 且 \(y = p_{\operatorname{A}}(x)\) ；我们有 \(y \in \operatorname{dom}(u) \cap \operatorname{E}_{\operatorname{A}}\) ，依据 \(c\) )。\(u(y)\) 关于 \(u\) 的谱测度是 \(|\operatorname{Id}_{\operatorname{Sp}(u)}|^2 \cdot \mu_y\) （loc. cit.）。由于 \(\mu_y\) 的支撑在 A 中，\(\mu_{u(y)}\) 亦然，故 \(u(y) \in \operatorname{E}_{\operatorname{A}}\) ，依据 \(a\) )。

推论. --- 设 \(\lambda \in \mathrm{Sp}(u)\) 。\(p_{\{\lambda\}}\) 的像是 \(u\) 关于 \(\lambda\) 的特征子空间。

设 \(x \in \operatorname{dom}(u)\) 。用 \(\mu\) （相应地，\(\nu\) ）表示元素 \(x\) 关于 \(u\) 的谱测度（相应地，\((x, u(x))\) 关于 \(u\) 的谱测度）。我们有 \(\nu = \operatorname{Id}_{\operatorname{Sp}(u)} \cdot \mu\) （IV, p. 275 的命题 5 的推论）。若 \(u(x) = \lambda x\) ，我们还有 \(\nu = \lambda \mu\) ，由此得等式 \(\operatorname{Id}_{\operatorname{Sp}(u)} \cdot \mu = \lambda \mu\) 。这蕴含 \(\mu\) 的支撑包含在 \(\{\lambda\}\) 中（参见 INT, V, p. 46, § 5, 第 3 目, 推论 2），因此 \(x\) 属于 \(p_{\{\lambda\}}\) 的像（命题 9, a)）。

反之，假设 \(x\) 属于 \(\mathrm{E}_{\lambda}\) ，即 \(p_{\{\lambda\}}\) 的像。那么 \(x\) 属于 \(\operatorname{dom}(u)\) ，且 \(u(x)\) 也属于 \(\mathrm{E}_{\lambda}\) （loc. cit., \(b\) )）。由于 \(\varphi_{\{\lambda\}} \cdot (\operatorname{Id}_{\operatorname{Sp}(u)} - \lambda) = 0\) ，我们有关系式 \(p_{\{\lambda\}} \circ (u - \lambda 1_{\mathrm{E}}) \subset 0\) （IV, p. 275 的命题 5, \(c\) )），由此 \(0 = p_{\{\lambda\}}(u(x)) - \lambda p_{\{\lambda\}}(x) = u(x) - \lambda x\) 。

命题 10. — 设 \(\lambda \in \mathbf{C}\) 。我们有 \(\lambda \in \mathrm{Sp}(u)\) 当且仅当：对 \(\lambda\) 在 \(\mathbf{C}\) 中的每个开邻域 V，u 关于 V 的谱投影 \(p_{\mathrm{V}}\) 非零。

若 \(\lambda \notin \mathrm{Sp}(u)\) ，则存在开邻域 \(V\) ，它包含 \(\lambda\) ，位于 \(\mathbf{C}\) 中，且与 \(\mathrm{Sp}(u)\) 不相交，于是 \(p_{\mathrm{V}} = p_{\varnothing} = 0\) 。

反之，假设存在 \(\lambda\) 在 \(\mathbf{C}\) 中的开邻域 V 使得 \(p_{V} = 0\) 。设 c > 0 使得以 \(\lambda\) 为中心、以 c 为半径的圆盘包含在 V 中。

设 \(x \in \mathrm{dom}(u)\) ，并设 \(\mu_x\) 是 \(x\) 关于 \(u\) 的谱测度。由于 \(\mu_x(\mathrm{V}) = \langle x | p_{\mathrm{V}}(x) \rangle = 0\) ，我们计算出

\[
\begin{array}{r l} & {\int_ {\mathbf {C}} | z - \lambda | ^ {2} d \mu_ {x} (z) = \int_ {\mathbf {C - V}} | z - \lambda | ^ {2} d \mu_ {x} (z)} \\ & {\qquad \geqslant c ^ {2} \int_ {\mathbf {C - V}} d \mu_ {x} (z) = c ^ {2} \int_ {\mathbf {C}} d \mu_ {x} (z) = c ^ {2} \| x \| ^ {2}.} \end{array}
\]

但是，另一方面，我们有

\[
\| u (x) - \lambda x \| ^ {2} = \int_ {\mathbf {C}} | z - \lambda | ^ {2} d \mu_ {x} (z) = \| u ^ {*} (x) - \overline {{\lambda}} x \| ^ {2}
\]

（p. 272 的公式 (8)），由此 \(\|u(x)-\lambda x\|\geqslant c\|x\|\) 且 \(\|u^{*}(x)-\overline{\lambda}x\|\geqslant c\|x\|\) 。于是 \(\lambda\) 属于 u 的预解集（IV, p. 248 的引理 5）。证明完毕。

推论. — 设 A 是 \(\mathbf{C}\) 中的开集，且 \(n \in \mathbf N\) 。若 A 含有 \(\mathrm{Sp}(u)\) 的 n 个元素，则 u 关于 A 的谱投影 \(p_{A}\) 的像的维数至少等于 n。

设 \(\lambda_1, \ldots, \lambda_n\) 是 \(u\) 的谱中属于 A 的互不相同的元素。存在两两不相交的开集族 \((\mathrm{V}_i)_{1 \leqslant i \leqslant n}\) （\(\mathbf{C}\) 中的开集），使得 \(\lambda_i \in \mathrm{V}_i\) 对 \(1 \leqslant i \leqslant n\) 成立。设 B 是诸集合 \(\mathrm{V}_i\) 的并。\(p_{\mathrm{A}}\) 的像包含谱投影 \(p_{\mathrm{B}}\) 的像；此外，由于 \(p_{\mathrm{B}}\) 是诸投影 \(p_{\mathrm{V}_i}\) 之和，且 \(p_{\mathrm{V}_i}\) 的像与 \(p_{\mathrm{V}_j}\) 的像正交（对所有 \(i \neq j\) ），结果由命题 10 得出。

